\documentclass[12pt,a4paper]{article}
\usepackage[margin=2cm]{geometry}
\usepackage[utf8]{inputenc}
\usepackage[T1]{fontenc}
\usepackage{titlesec}

\title{Conspecte Biblice: Geneza, Ieșirea, Levitic, Numeri, Iosua}
\author{Ștefan Darius Țacu}
\date{\today}

\begin{document}

\maketitle

\section{Geneza}

\subsection{Capitolul 1}

Dumnezeu a făcut lumea în 6 zile.

Dumnezeu S-a sfătuit să facă pe om dupa chipul și asemănarea Sa și
a făcut pe om după chipul Său.

\subsection{Capitolul 2}

Dumnezeu poruncește omului să nu mănânce din pomul cunoștinței binelui și răului.

Dumnezeu aduce somn greu asupra lui Adam și o face pe femeie din coasta lui.

\subsection{Capitolul 3}

Din dorința de a fi ca Dumnezeu, Adam și Eva calcă porunca.

Adam și Eva sunt izgoniți din Rai.

\subsection{Capitolul 4}

Cain îl omoară pe fratele său, Abel.

Pentru fapta sa, Dumnezeul pune un semn distinctiv asupra lui Cain.

\subsection{Capitolul 5}

Genealogia de la Adam până la Noe este: Adam, Set, Enos, Cainan, Maleleil,
Iared, Enoh, Metuselem, Lameh, Noe.

Lameh îl naște pe Noe la vârsta de 188 de ani.

\subsection{Capitolul 6}

Lui Dumnezeu îi pare rău că a făcut pe om și decide să-l piardă de pe fața pământului.

Dumnezeui poruncește lui Noe să construiască o corabie pentru a rămâne el și familia sa în viață.

\subsection{Capitolul 7}

Noe aduce în corabie 7 perechi de animale curate și o pereche de animale necurate.

Potopul a omorât toți oamenii și toate vietățile cele de pe uscat.

\subsection{Capitolul 8}

Noe împreună cu familia sa și toate vietățile ies din corabie după ce au secat apele de pe pământ.

Noe aduce jertfă lui Dumnezeu, iar Dumnezeu făgăduiește că nu va mai nimici pământul pentru faptele omului.

\section{Ieșirea}

\subsection{Capitolul 1}
% Placeholder pentru conspect: Ieșirea 1
% Adaugă aici conspectul pentru Ieșirea capitolul 1.

\subsection{Capitolul 2}
% Placeholder pentru conspect: Ieșirea 2
% Adaugă aici conspectul pentru Ieșirea capitolul 2.

% Continuă cu capitolele 3 până la 40.

\section{Levitic}

\subsection{Capitolul 1}
% Placeholder pentru conspect: Levitic 1
% Adaugă aici conspectul pentru Levitic capitolul 1.

\subsection{Capitolul 2}
% Placeholder pentru conspect: Levitic 2
% Adaugă aici conspectul pentru Levitic capitolul 2.

% Continuă cu capitolele 3 până la 27.

\section{Numeri}

\subsection{Capitolul 1}
% Placeholder pentru conspect: Numeri 1
% Adaugă aici conspectul pentru Numeri capitolul 1.

\subsection{Capitolul 2}
% Placeholder pentru conspect: Numeri 2
% Adaugă aici conspectul pentru Numeri capitolul 2.

% Continuă cu capitolele 3 până la 36.

\section{Iosua}

\subsection{Capitolul 1}
% Placeholder pentru conspect: Iosua 1
% Adaugă aici conspectul pentru Iosua capitolul 1.

\subsection{Capitolul 2}
% Placeholder pentru conspect: Iosua 2
% Adaugă aici conspectul pentru Iosua capitolul 2.

% Continuă cu capitolele 3 până la 24.

\end{document}
